\documentclass[10pt,letterpaper]{article}
\usepackage{cornell-notes}

\title{Clause 03: Terms, Definitions, and Symbols}
\author{}
\date{}

\setDocTitle{Clause 03: Terms, Definitions, and Symbols}
\setDocSubtitle{{C 2024}}
\setDocOwner{{C 2024 Cornell Notes}}
\setDocDate{{\today}}
\setCornellCollection{{C 2024 Cornell Notes}}
\setCornellUnitType{{Clause}}
\setCornellUnitNumber{03}
\setCornellUnitTitle{Terms, Definitions, and Symbols}
\hypersetup{pdftitle={Clause 03: Terms, Definitions, and Symbols},pdfsubject={C 2024 study notes},pdfauthor={}}

\begin{document}
\maketitle



\begin{CornellOverviewBox}{Clause Orientation}
This clause supplies the vocabulary needed to classify program behavior, storage, values, diagnostics, concurrency, and bounds failures. The distinctions are operational: confusing undefined, unspecified, implementation-defined, and locale-specific behavior changes what a programmer may assume and what an implementation must document.
\end{CornellOverviewBox}

\section{Learning Objectives}
\begin{itemize}
  \item Define access, object, value, byte, bit, and alignment in the specification sense.
  \item Distinguish the four principal behavior classifications.
  \item Differentiate arguments from parameters and constraints from runtime-constraints.
  \item Explain memory locations and their importance to concurrent access.
  \item Relate unspecified values, indeterminate representations, non-value representations, and traps.
  \item Classify bounded and critical undefined behavior using the out-of-bounds-store boundary.
  \item Apply ceiling, floor, and wraparound notation correctly.
\end{itemize}

\section{Normative Structure Map}
\begin{CornellNotesTable}
\CornellNoteRow{Core definitions}{Access, object, alignment, argument, and parameter all have precise semantic roles.}
\CornellNoteRow{Behavior taxonomy}{Implementation-defined, locale-specific, unspecified, and undefined behavior are distinct categories.}
\CornellNoteRow{Representation vocabulary}{Bits, bytes, characters, and memory locations define the behavioral meaning of values.}
\end{CornellNotesTable}

\section{Objects, Access, Alignment, and Calls}
\begin{CornellNotesTable}
\CornellNoteRow{What counts as an access?}{At execution time, reading or modifying an object is an access. Storing the same value still counts as modification; an unevaluated expression does not access an object.}
\CornellNoteRow{What is an object?}{A region of execution-environment storage whose contents can represent values. A use of the object interprets those contents through a type.}
\CornellNoteRow{What is alignment?}{A placement requirement that constrains addresses of objects of a type to particular storage boundaries.}
\CornellNoteRow{Argument or parameter?}{An argument is an expression supplied to a function call or a preprocessing-token sequence supplied to a function-like macro. A parameter is the declared function object or macro identifier that receives it.}
\CornellNoteRow{What does arithmetic negation mean?}{It produces the negative of a number. For an integer it corresponds to subtraction from zero; for floating types it changes the sign subject to the applicable arithmetic rules.}
\end{CornellNotesTable}

\begin{CornellSummaryBox}{Section Summary}
The vocabulary separates runtime actions and storage from the syntactic roles used to pass information into functions and macros.
\end{CornellSummaryBox}



\section{Behavior Taxonomy}
\begin{CornellNotesTable}
\CornellNoteRow{What is behavior?}{The externally apparent action or result associated with translation or execution.}
\CornellNoteRow{What is implementation-defined behavior?}{It is an unspecified choice for which each implementation must document how the choice is made.}
\CornellNoteRow{What is locale-specific behavior?}{It depends on local cultural, national, or language conventions and must be documented by the implementation.}
\CornellNoteRow{What is unspecified behavior?}{The source permits two or more possibilities, or an unspecified value determines the result, without requiring the implementation to document which possibility occurs in a particular instance.}
\CornellNoteRow{What is undefined behavior?}{The source imposes no requirements after the relevant erroneous or nonportable construct or data is encountered. A diagnostic is not inherently required unless another rule, such as a violated constraint, requires one.}
\CornellNoteRow{What remains true before undefined behavior?}{Observable behavior that occurs before the operation producing undefined behavior remains governed by the specification; later consequences can follow from the concrete behavior chosen when the error is encountered.}
\end{CornellNotesTable}

\begin{CornellSummaryBox}{Section Summary}
Documentation is the key discriminator: implementation-defined and locale-specific choices must be documented; unspecified choices need not be; undefined behavior supplies no general behavioral guarantee.
\end{CornellSummaryBox}



\section{Bits, Bytes, and Characters}
\begin{CornellNotesTable}
\CornellNoteRow{How do bit and byte differ?}{A bit can hold one of two values but need not be individually addressable. A byte is an addressable storage unit large enough to hold any basic execution character.}
\CornellNoteRow{What are low- and high-order bits?}{They are respectively the least and most significant bits within a byte.}
\CornellNoteRow{What is an abstract character?}{An element used to organize, control, or represent data, independent of a particular encoding.}
\CornellNoteRow{Single-byte, multibyte, or wide character?}{A single-byte character fits in one byte; a multibyte character is a byte sequence in an extended character set; a wide character is a value of \texttt{wchar\_t} capable of representing any character in the current locale.}
\end{CornellNotesTable}

\begin{CornellSummaryBox}{Section Summary}
Character concepts distinguish abstract repertoire members from the byte sequences or wide values used to encode them.
\end{CornellSummaryBox}



\section{Constraints, Diagnostics, and Implementations}
\begin{CornellNotesTable}
\CornellNoteRow{What is a constraint?}{A syntactic or semantic restriction used to interpret a language element. Violating a constraint triggers the diagnostic rule in the translation environment.}
\CornellNoteRow{What is a diagnostic message?}{A message belonging to an implementation-defined subset of the implementation output designated as diagnostics.}
\CornellNoteRow{What is recommended practice?}{Strong guidance consistent with the specification intent that may be impractical for some implementations and is therefore not an unconditional requirement.}
\CornellNoteRow{What is a runtime-constraint?}{A requirement on a program when calling a library function. Despite the name, it is distinct from a language constraint and is not necessarily diagnosed during translation.}
\CornellNoteRow{What is an implementation?}{A particular translator and support environment operating under particular controls for a particular execution environment.}
\CornellNoteRow{What is an implementation limit?}{A restriction the implementation imposes on programs, subject to required minima and documentation obligations.}
\CornellNoteRow{What is a forward reference?}{A pointer to a later subclause containing additional relevant information; it guides reading but does not replace the current rule.}
\end{CornellNotesTable}

\begin{CornellSummaryBox}{Section Summary}
Constraint, runtime-constraint, diagnostic, recommendation, and implementation limit are different rule categories and must not be collapsed into a generic notion of error.
\end{CornellSummaryBox}



\section{Memory Locations and Concurrent Access}
\begin{CornellNotesTable}
\CornellNoteRow{What is a memory location?}{Either a scalar object or a maximal adjacent sequence of nonzero-width bit-fields. It is the granularity used when reasoning about potentially interfering accesses.}
\CornellNoteRow{Why do adjacent bit-fields matter?}{Two non-atomic bit-fields within the same memory location cannot be modified concurrently without a data-race risk. A zero-width bit-field, non-bit-field member, or certain nesting boundaries can separate memory locations.}
\CornellNoteRow{What concurrency inference is safe?}{Separate memory locations can be accessed and updated independently with respect to the memory-location definition, though other synchronization and lifetime rules still apply.}
\end{CornellNotesTable}

\begin{CornellSummaryBox}{Section Summary}
Memory location is a semantic unit for concurrency, not necessarily a byte, machine word, source declaration, or cache line.
\end{CornellSummaryBox}



\section{Values, Representations, and Traps}
\begin{CornellNotesTable}
\CornellNoteRow{What is a value?}{The meaning obtained when an object representation is interpreted as a particular type.}
\CornellNoteRow{What is an implementation-defined value?}{An unspecified value for which the implementation documents its selection.}
\CornellNoteRow{What is an unspecified value?}{Any valid value of the relevant type, with no requirement that the implementation document which one is selected in an instance.}
\CornellNoteRow{What is an indeterminate representation?}{An object representation that denotes either an unspecified value or a non-value representation.}
\CornellNoteRow{What is a non-value representation?}{A representation that does not encode a value of the object type.}
\CornellNoteRow{What does it mean to perform a trap?}{Execution is interrupted so that no further operations occur. Fetching a non-value representation can permit, but does not always require, a trap.}
\end{CornellNotesTable}

\begin{CornellSummaryBox}{Section Summary}
Always reason in the order representation, type interpretation, value status, and permitted behavior; indeterminate does not mean one single runtime outcome.
\end{CornellSummaryBox}



\section{Bounds and Undefined-Behavior Severity}
\begin{CornellNotesTable}
\CornellNoteRow{What is an out-of-bounds store?}{An attempted runtime access that would modify bytes outside the bounds permitted by the specification; for a volatile object, the defined classification also accounts for fetching outside permitted bounds.}
\CornellNoteRow{What is bounded undefined behavior?}{Undefined behavior that does not perform an out-of-bounds store. It can still trap, yield unspecified values, or leave written representations indeterminate.}
\CornellNoteRow{What is critical undefined behavior?}{Undefined behavior that is not bounded. It can perform an out-of-bounds store or perform a trap.}
\CornellNoteRow{What is wraparound?}{Reduction modulo \texttt{2\textasciicircum{}N}, where \texttt{N} is the width of the resulting type. It is a defined mathematical process only where the governing type and operation rules specify it.}
\CornellNoteRow{What do ceiling and floor denote?}{Ceiling selects the least integer not less than a number; floor selects the greatest integer not greater than it.}
\end{CornellNotesTable}

\begin{CornellSummaryBox}{Section Summary}
The bounded-versus-critical distinction classifies undefined behavior by whether an out-of-bounds store can occur; it does not make bounded undefined behavior correct or portable.
\end{CornellSummaryBox}



\section{Programmer and Implementation Checklist}
\begin{itemize}
  \item[$\square$] Classify each disputed result as specified, unspecified, implementation-defined, locale-specific, or undefined.
  \item[$\square$] Check whether the implementation owes documentation for the chosen behavior or value.
  \item[$\square$] Do not equate a runtime-constraint with a translation-time constraint.
  \item[$\square$] Determine whether an expression is evaluated before claiming it accesses an object.
  \item[$\square$] Reason about concurrency at the memory-location granularity, especially for adjacent bit-fields.
  \item[$\square$] Separate object representation from the value obtained through a type.
  \item[$\square$] Identify whether an indeterminate representation could be a non-value representation.
  \item[$\square$] Classify undefined behavior as bounded or critical only after testing the out-of-bounds-store condition.
  \item[$\square$] Do not infer that wraparound applies to every integer computation.
  \item[$\square$] Preserve the difference between a recommendation and a mandatory requirement.
\end{itemize}

\begin{CornellSummaryBox}{Clause Synthesis}
Clause 3 is the vocabulary engine for the rest of the material. Its behavior taxonomy determines what programs may assume and what implementations must document. Its storage and representation definitions determine how values, concurrency, and bounds are analyzed. Precise use of these terms prevents category errors such as treating unspecified behavior as undefined, treating a runtime-constraint as a language constraint, or treating bounded undefined behavior as safe.
\end{CornellSummaryBox}



\clearpage
\section{Self-Test Questions}
\begin{enumerate}
  \item Does storing the same value back into an object count as modification?
  \item What distinguishes implementation-defined behavior from unspecified behavior?
  \item Does undefined behavior always require a diagnostic?
  \item Why is a byte not defined as eight bits?
  \item How does a parameter differ from an argument?
  \item Can a runtime-constraint violation be assumed to produce a translation diagnostic?
  \item Two adjacent non-atomic bit-fields are updated by different threads. What must be checked first?
  \item Does every indeterminate representation denote an unspecified value?
  \item What condition separates bounded from critical undefined behavior?
  \item Why is bounded undefined behavior still unacceptable in portable code?
  \item An implementation documents the result of a right-shift choice. Which category is indicated?
  \item A library call violates a stated requirement on its arguments at runtime. Which vocabulary should be checked?
\end{enumerate}

\clearpage
\section{Answer Key}
\begin{enumerate}
  \item Yes. Modification includes storing a value even when it equals the previous value, so the action is an access under 3.1.

  \item Both permit a choice, but implementation-defined behavior requires the implementation to document how the choice is made.

  \item No. A diagnostic is required only when another rule requires one, such as a syntax-rule or constraint violation.

  \item The definition is functional: it is the addressable unit large enough for a basic execution character. Its bit width is constrained elsewhere rather than fixed by the term itself.

  \item A parameter is declared by a function or macro definition; an argument is the expression or token sequence supplied at a call or invocation.

  \item No. A runtime-constraint is distinct from a language constraint and need not be diagnosed during translation.

  \item Determine whether the fields occupy the same memory location under 3.17. If they do, concurrent conflicting access requires synchronization to avoid a data race.

  \item No. It can instead be a non-value representation, which may permit a trap when fetched.

  \item Bounded undefined behavior does not perform an out-of-bounds store. Critical undefined behavior is undefined behavior that is not bounded.

  \item The source still imposes no ordinary behavioral requirements; a trap, unspecified values, or indeterminate stored representations can occur.

  \item If the source designates that choice as implementation-defined, the documented selection is implementation-defined behavior.

  \item Check whether the provision is a library precondition or explicitly a runtime-constraint; do not automatically classify it as a translation-time constraint.
\end{enumerate}

\section{Key-Term Recap}
\begin{CornellNotesTable}
\toprule
Term & Working meaning \\
\midrule
\endfirsthead
\toprule
Term & Working meaning (continued) \\
\midrule
\endhead
Access & A runtime read or modification of an object. \\
Object & A region of storage whose contents can represent values. \\
Memory location & A scalar object or maximal adjacent sequence of nonzero-width bit-fields. \\
Implementation-defined behavior & An unspecified choice that the implementation documents. \\
Locale-specific behavior & Documented behavior dependent on local cultural or language conventions. \\
Unspecified behavior & One of multiple permitted possibilities with no required per-instance documentation. \\
Undefined behavior & Behavior for which the source imposes no requirements. \\
Constraint & A syntactic or semantic restriction governing a language element. \\
Runtime-constraint & A requirement imposed on a program when it calls a library function. \\
Indeterminate representation & A representation denoting either an unspecified value or a non-value representation. \\
Non-value representation & An object representation that encodes no value of the object type. \\
Trap & Termination of execution such that no further operations occur. \\
Wraparound & Reduction modulo a power of two determined by the result type width. \\
Out-of-bounds store & An attempted access that would modify bytes outside permitted bounds. \\
Bounded undefined behavior & Undefined behavior that performs no out-of-bounds store. \\
Critical undefined behavior & Undefined behavior that is not bounded. \\
\bottomrule
\end{CornellNotesTable}

\begin{CornellSummaryBox}{One-Sentence Takeaway}
Precise classification of behavior, storage, and representation is the foundation for every later conclusion about conformance, portability, and program correctness.
\end{CornellSummaryBox}
\end{document}
